\documentclass[UTF8,12pt,a4paper,fontset=fandol]{ctexbook}
\usepackage[margin=2.5cm]{geometry}
\usepackage{setspace}
\usepackage{graphicx}
\usepackage{longtable}
\usepackage{booktabs}
\usepackage{array}
\usepackage[hidelinks]{hyperref}
\setstretch{1.35}
\setlength{\parindent}{2em}

\title{北京理工大学本科生毕业设计（论文）LaTeX 备用模板（由 DOCX 转换）}
\author{}
\date{}

\begin{document}
\maketitle
\tableofcontents
\clearpage

% 说明：本文件为从 Word 模板自动提取的 LaTeX 备用稿，需按实际论文内容进一步精修版式。
% 原始模板文件已备份于同级目录。

本科生毕业设计(论文)

北京理工大学本科生毕业设计（论文）题目

\begin{center}\small The Subject of Undergraduate Graduation Project (Thesis) of Beijing Institute of Technology\end{center}

\begin{longtable}{p{0.18\textwidth}p{0.18\textwidth}}
\toprule
学    院： &  \\
\midrule
专    业： &  \\
班    级： &  \\
学生姓名： &  \\
学    号： &  \\
指导教师： &  \\
\bottomrule
\end{longtable}

20XX 年   月   日

本人郑重声明：所呈交的毕业设计（论文），是本人在指导老师的指导下独立进行研究所取得的成果。除文中已经注明引用的内容外，本文不包含任何其他个人或集体已经发表或撰写过的研究成果。对本文的研究做出重要贡献的个人和集体，均已在文中以明确方式标明。

特此申明。

本人签名：               日 期：     年   月   日

关于使用授权的声明

本人完全了解北京理工大学有关保管、使用毕业设计（论文）的规定，其中包括：①学校有权保管、并向有关部门送交本毕业设计（论文）的原件与复印件；②学校可以采用影印、缩印或其它复制手段复制并保存本毕业设计（论文）；③学校可允许本毕业设计（论文）被查阅或借阅；④学校可以学术交流为目的,复制赠送和交换本毕业设计（论文）；⑤学校可以公布本毕业设计（论文）的全部或部分内容。

本人签名：               日 期：     年   月   日

指导老师签名：               日 期：     年   月   日

本文……。

摘要正文选用模板中的样式所定义的“正文”，每段落首行缩进2个字符；或者手动设置成每段落首行缩进2个汉字，字体：宋体，字号：小四，行距：固定值22磅，间距：段前、段后均为0行。【阅后删除此段】

摘要是一篇具有独立性和完整性的短文，应概括而扼要地反映出本论文的主要内容。包括研究目的、研究方法、研究结果和结论等，特别要突出研究结果和结论。中文摘要力求语言精炼准确，本科生毕业设计（论文）摘要建议300-500字。摘要中不可出现参考文献、图、表、化学结构式、非公知公用的符号和术语。英文摘要与中文摘要的内容应一致。【阅后删除此段】

\chapter*{Abstract}
\addcontentsline{toc}{chapter}{Abstract}

In order to study……

摘　要I

AbstractII

第1章 一级题目1

1.1 二级题目1

1.1.1 三级题目1

结　论3

参考文献4

附　录6

致　谢7

\chapter{第1章 一级题目}

\section{1.1 二级题目}

\subsection{1.1.1 三级题目}

正文……

正文部分：宋体、小四；正文行距：22磅；间距段前段后均为0行。【阅后删除此段】

图、表居中，图注标在图下方，表头标在表上方，宋体、五号、居中，1.25倍行距，间距段前段后均为0行，图表与上下文之间各空一行。【阅后删除此段】

图-示例：【阅后删除此段】

\begin{center}\small 图1-1 标题序号\end{center}

表-示例：【阅后删除此段】

\begin{center}\small 表1-1 统计表\end{center}

\begin{longtable}{p{0.18\textwidth}p{0.18\textwidth}p{0.18\textwidth}p{0.18\textwidth}p{0.18\textwidth}}
\toprule
项目 & 产量 & 销量 & 产值 & 比重 \\
\midrule
手机 & 1000 & 10000 & 500 & 50\% \\
计算机 & 5500 & 5000 & 220 & 22\% \\
笔记本电脑 & 1100 & 1000 & 280 & 28\% \\
合计 & 17600 & 16000 & 1000 & 100\% \\
\bottomrule
\end{longtable}

公式标注应于该公式所在行的最右侧。对于较长的公式只可在符号处（+、-、*、/、≤≥等）转行。在文中引用公式时，在标号前加“式”，如式（1-2）。【阅后删除此段】

公式-示例：【阅后删除此段】

(1-1)

本文结论……。

结论作为毕业设计（论文）正文的最后部分单独排写，但不加章号。结论是对整个论文主要结果的总结。在结论中应明确指出本研究的创新点，对其应用前景和社会、经济价值等加以预测和评价，并指出今后进一步在本研究方向进行研究工作的展望与设想。结论部分的撰写应简明扼要，突出创新性。【阅后删除此段】

结论正文样式与文章正文相同：宋体、小四；行距：22磅；间距段前段后均为0行。【阅后删除此段】

参考文献书写规范

参考国家标准《信息与文献参考文献著录规则》【GB/T 7714—2015】，参考文献书写规范如下：

1. 文献类型和标识代码

普通图书：M     会议录：C       汇编：G          报纸：N

期刊：J        学位论文：D       报告：R          标准：S

专利：P        数据库：DB     计算机程序：CP   电子公告：EB

档案：A         舆图：CM       数据集：DS        其他：Z

2. 不同类别文献书写规范要求

期刊

[序号] 主要责任者. 文献题名[J]. 刊名, 出版年份, 卷号(期号): 起止页码.

[1] 余雄庆. 飞机总体多学科设计优化的现状与发展方向[J]. 南京航空航天大学学报, 2008, 40(4): 417-426.

[2] Hajela P, Bloebaumj C L, Sobieszczanski-Sobieski J. Application of Global Sensitivity Equations in Multidisciplinary Aircraft Synthesis[J]. Journal of Aircraft, 1990, 27(12): 1002-110.

普通图书

[序号] 主要责任者. 文献题名[M]. 出版地: 出版者, 出版年: 起止页码.

[3] 张伯伟. 全唐五代诗格会考[M]. 南京: 江苏古籍出版社, 2002: 288.

[4] O’BRIEN J A. Introduction to information systems[M]. 7th ed. Burr Ridge, III: Irwin, 1994.

会议论文集

[序号] 主要责任者．题名:其他题名信息[C]. 出版地: 出版者, 出版年.

[5] 雷光春. 综合湿地管理: 综合湿地管理国际研讨会论文集[C]. 北京: 海洋出版社, 2012.

专著中析出的文献

[序号] 析出文献主要责任者. 析出题名[M]//专著主要责任者. 专著题名. 出版地: 出版者, 出版年: 起止页码.

[6] 白书农. 植物开花研究[M]//李承森. 植物科学进展. 北京: 高等教育出版社, 1998: 146-163.

学位论文

[序号] 主要责任者. 文献题名[D]. 保存地: 保存单位, 年份.

[7] 张和生. 嵌入式单片机系统设计[D]. 北京: 北京理工大学, 1998.

[8] Sobieski I P. Multidisciplinary Design Using Collaborative Optimization[D]. United States -- California: Stanford University, 1998.

报告

[序号] 主要责任者. 文献题名[R]. 报告地: 报告会主办单位, 年份.

[9] 冯西桥. 核反应堆压力容器的LBB分析[R]. 北京: 清华大学核能技术设计研究院, 1997.

[10] Sobieszczanski-Sobieski J. Optimization by Decomposition: A Step from Hierarchic to Non-Hierarchic Systems[R]. NASA CP-3031, 1989.

专利文献

[序号] 专利所有者. 专利题名:专利号[P]. 公告日期或公开日期[引用日期]. 获取和访问路径. 数字对象唯一标识符.

[11] 姜锡洲. 一种温热外敷药制备方案: 881056078 [P]. 1983-08-12.

国际、国家标准

[序号] 主要责任人. 题名: 其他题名信息[S]. 出版地: 出版者, 出版年: 引文页码.

[12] 全国信息与文献标准化技术委员会. 文献著录: 第4部分 非书资料: GB/T 3792.4-2009[S].  北京: 中国标准出版社, 2010: 3.

报纸文章

[序号] 主要责任者. 文献题名[N]. 报纸名, 年(期): 页码.

[13] 谢希德. 创造学习的思路[N]. 人民日报, 1998-12-25(10).

电子文献

[序号] 主要责任者. 电子文献题名[文献类型/载体类型]. (发表或更新日期) [引用日期]. 获取和访问路径. 数字对象唯一标识符.

[14] 姚伯元. 毕业设计(论文)规范化管理与培养学生综合素质[EB/OL]. [2005-02-02]. 中国高等教育网教学研究.

关于参考文献的未尽事项可参考国家标准《信息与文献参考文献著录规则》（GB/T 7714—2015）

附录相关内容…

附录是毕业设计（论文）主体的补充项目，为了体现整篇文章的完整性，写入正文又可能有损于论文的条理性、逻辑性和精炼性，这些材料可以写入附录段，但对于每一篇文章并不是必须的。附录依次用大写正体英文字母A、B、C……编序号，如附录A、附录B。【阅后删除此段】

附录正文样式与文章正文相同：宋体、小四；行距：22磅；间距段前段后均为0行。【阅后删除此段】

值此论文完成之际，首先向我的导师……

致谢正文样式与文章正文相同：宋体、小四；行距：22磅；间距段前段后均为0行。【阅后删除此段】

\end{document}